% !TeX root = ../sections/02_exercises.tex
\subsection{2-A-3.}
\begin{exercise}
	下に有界な単調減少数列
	\[
	\{a_n\}_{n\in\mathbb{N}}
	\]
	は
	\[
	\inf_{n\in\mathbb{N}} a_n
	\]
	に収束することを証明せよ。
\end{exercise}

\begin{background}
	単調減少数列とは，任意の \(n\in\mathbb{N}\) に対して
	\[
	a_{n+1}\leq a_n
	\]
	が成り立つ数列である。
	
	また，下に有界であるとは，ある実数 \(m\) が存在して，
	任意の \(n\in\mathbb{N}\) に対して
	\[
	m\leq a_n
	\]
	が成り立つことである。
	
	このとき集合
	\[
	A=\{a_n\mid n\in\mathbb{N}\}
	\]
	は下に有界なので，下限
	\[
	\alpha=\inf A
	\]
	が存在する。
	
	直感的には，単調減少数列は上から下へ下がっていくが，
	下限 \(\alpha\) より下には行けない。
	したがって，数列は \(\alpha\) に近づく。
\end{background}

\begin{strategy}
	\[
	\alpha=\inf_{n\in\mathbb{N}} a_n
	\]
	とおく。
	示すべきことは
	\[
	\lim_{n\to\infty}a_n=\alpha
	\]
	である。
	
	任意に \(\varepsilon>0\) を取る。
	\(\alpha=\inf A\) であるから，下限の特徴づけより，
	ある \(N\in\mathbb{N}\) が存在して
	\[
	a_N<\alpha+\varepsilon
	\]
	となる。
	
	一方，\(\alpha\) は下界なので，任意の \(n\) に対して
	\[
	\alpha\leq a_n
	\]
	である。
	
	さらに，数列は単調減少なので，\(n\geq N\) ならば
	\[
	a_n\leq a_N
	\]
	である。
	したがって，\(n\geq N\) に対して
	\[
	\alpha\leq a_n\leq a_N<\alpha+\varepsilon
	\]
	が成り立つ。
	
	よって
	\[
	0\leq a_n-\alpha<\varepsilon
	\]
	であり，
	\[
	|a_n-\alpha|<\varepsilon
	\]
	となる。
	したがって
	\[
	a_n\to\alpha
	\]
	である。
\end{strategy}